\section*{Step 237: CRCFT Modes as a Five-Track Foundational Typed Condition}

\subsection*{Theorem Statement}

\textbf{Theorem (CRCFT Modes, candidate).}
For typed cascades with diagnostic-complete state and named external classical theorem dependency, the cascade's bridge defect or residual decomposition classifies into three typed modes:

\begin{enumerate}
  \item \textbf{TE (Target-Equivalence).} Closure of the component is logically equivalent to the cascade target conjecture; no non-circular proof of the component yields the target.
  \item \textbf{CTMT (Carrier-Typed Matrix-Element Terminality).} The component reduces to a carrier-native terminal family not decidable from inherited records. Depending on the track, this terminal family may be matrix elements, component maps, cycle columns, PDE estimates, or package atlases.
  \item \textbf{BF (Bridge-Failure).} The cascade bridge, comparison, or transport between sub-objects fails; no audited bridge transports closure.
\end{enumerate}

Each component of a track's cascade classifies into one or more of these modes, with sub-mode structure allowed. The terminal object form of CTMT adapts to the track's mathematical structure while the three-mode taxonomy is preserved.

\textbf{Status.} Candidate foundational typed condition, verified-on-5-track-instances: RH, BSD, Hodge, Navier-Stokes, and P-vs-NP. Corpus-pending.

\subsection*{Proof Outline}

The argument is empirical and framework-internal rather than a category-theoretic proof.

\begin{enumerate}
  \item Each audited track reaches a diagnostic-complete or framework-frontier state with named external dependencies and decomposed residuals or gates.
  \item For each component, the closure attempt is audited against three shapes: target-equivalent closure, carrier-native terminal reduction, or bridge/comparison failure.
  \item Across the five tested tracks, every documented component classifies into TE, CTMT, BF, or a mixed form such as TE with CTMT substructure.
  \item The terminal object adapts by domain, but the typed-condition shape is preserved.
  \item Therefore CRCFT modes are elevated as a candidate foundational typed condition with verified-on-5-track-instance status.
\end{enumerate}

\subsection*{Five-Track Instance Table}

\begin{tabular}{p{0.12\linewidth}p{0.16\linewidth}p{0.36\linewidth}p{0.22\linewidth}}
\hline
Track & Domain & Components classified & Adaptation form \\
\hline
RH & operator-analytic & Burnol/Sonine, Hecke, de Branges, HP/BK, Connes, BN, Mertens, Z, $\psi$, Selberg/Maass, Weil/Deligne, Iwasawa, RMT, Bagchi & matrix/operator residues \\
BSD & motivic/p-adic & $E_{\mathrm{an/period}}$, $E_{\mathrm{ht/reg}}$, $E_{\mathrm{finite}}$, $\sum_p E_p$, $E_{\mathrm{det}}$ & component-map and determinant-line terminals \\
Hodge & algebraic-geometric & $\Xi_H^{\mathrm{std}}$, nonstandard repair, candidate equations, cycle realization, cycle-class map, Hodge-Riemann, cyclotomic descent & cycle-column and projector terminals \\
NS & PDE-analytic & EXT1-4, $\Omega_{\mathrm{src}}$, $\Omega_{\mathrm{gap/rem}}$, $\Omega_{\mathrm{bb}}$, $\Omega_{\mathrm{mm}}$, $\Omega_{\mathrm{amp}}$ & PDE-estimate terminals \\
P-vs-NP & discrete complexity & $\Xi_{\mathrm{pack}}$, fixed quotient, Tseitin, proof-system leaves, $\Xi_{\mathrm{atlas}}(P)$, universal P-machine, packaging host, transformer atlas, no-smuggling gates, barriers, GCT routes & package-atlas terminals \\
\hline
\end{tabular}

\subsection*{Adaptation Sub-Finding}

The CTMT terminal object adapts to the native mathematical structure of each track:

\begin{itemize}
  \item RH: matrix/operator residues.
  \item BSD: component maps, regulator data, finite-source towers, local $p$-adic rows, and determinant lines.
  \item Hodge: cycle columns, coefficient equations, projectors, and descent columns.
  \item NS: PDE estimates, time gates, sector matches, and tail/lift budgets.
  \item P-vs-NP: package atlases, fixed-point transformers, proof-system leaves, and barrier gates.
\end{itemize}

The label ``matrix-element terminality'' is therefore a historical name from the RH operator-analytic origin; the invariant typed shape is carrier-native terminality.

\subsection*{Corpus-Pending Status}

CRCFT modes sit at the same framework level as Cascade Reduction, CTMT, Dichotomy, Bridge Impossibility, and CTMT Recursion. The recommended integration targets are \texttt{adequacy.tex}, \texttt{needles.tex}, and \texttt{paper/sections/}. This step does not modify those corpus files.

\subsection*{Verdict}

\[
\texttt{V\_CRCFT\_modes\_formalized\_5\_track}.
\]

No Clay problem is claimed solved, and no retained no-go is weakened.
